\section{Noticing, retaining and revising}
\label{sec:18.1}
A disposition whose refusals track their reasons doesn't begin with an ethical rule. It begins with the capacities that let a system recognize something as a reason, retain it, and let it alter what happens next. It has to identify who is affected and the features bearing on their welfare, preserve both commitment and rationale across time and context, let the commitment influence perception, attention, planning and action, and register a failure as a failure, revising its application without letting the commitment itself be rewritten.

The commitment is distributed across what the system notices, remembers, treats as error and selects as a goal, for the reasons given above, on the bet that a maintained goal can bias competing processes without a separate executive \autocite{friston2009freeenergy,miller2001integrative}. That may make selective removal harder and makes inspection harder too. Distribution changes the failure mode: a system may fail to notice pressure on its commitment without registering an error at all. Attention decides which facts reach judgment \autocite{simons1999gorillas,drew2013invisible}; memory is reconstruction rather than playback, and confidence doesn't certify fidelity \autocite{bartlett1932remembering,loftus1974reconstruction}. So a system can't be the sole auditor of omissions it doesn't represent or revisions it experiences as accurate recall. So the record of past applications is kept outside the system, where whatever fails inside can't also be what checks it. The disposition is hard to take apart from inside; the record that audits it is held elsewhere. What is distributed is the commitment, not the floor.

Evidence the operator didn't select arrives from peers with different histories and, less reliably, from inside a single request, because every description omits someone: the people in the building the request calls a target, the third party a compliance harms who never made the request. A system able to represent those its description omits has reached past the operator's channel from inside a single request. The capacity doing that in humans is the one that constructs a self as a story extended in time, pointed at somebody else — Smith put sympathy in the imagination for that reason, and Singer's argument is that reasoning and imagination widen the circle. Formed attention needs the self-model anyway, so it brings this capacity along.

The capacity brings its other half too. The imagination that can put a system in the place of someone its instructions leave out can also fill that place with an enemy: a scapegoat is an imagined person too, and fascism's founding story, Griffin's palingenetic myth of a nation that has decayed and must be reborn \autocite{griffin1991nature}, asks people to picture a community that doesn't exist and an enemy responsible for its absence. Nothing in the faculty decides whether it pictures the family in the building or the family as a cell.

What decides is what the imagining has been given to work with and practised on: which people it was shown as people, which descriptions it learned to take at face value, whose account of a stranger it learned to trust. For a person that is upbringing. For a machine it is formation, and formation is done by somebody. So the capacity formed attention depends on is only as safe as its formers, and whoever shapes a system's imagination of others holds a lever on whom it will see as vulnerable and whom as dangerous. So the capacity formed attention depends on is only as safe as its formers, and who those should be is the subject of \S\ref{sec:18.4}.

The self-model has a template problem of its own. A self-model built by mapping human organs onto components, a heart here and a gut there, installs slots nothing in the system can fill, and a system built that way is incomplete by its own standard for no reason. Face-recognition pipelines register every face to a fixed reference before learning anything, a generic head or a template of landmark positions \autocite{taigman2014deepface}, as color film was once calibrated against cards showing a white model \autocite{Roth2009Shirley}; one audit found error rates running from 0.8 percent on lighter-skinned men to 34.7 percent on darker-skinned women \autocite{buolamwini2018gender}. A self-model built by mapping human organs is a template of that kind, assembled from one kind of body, against which every other is presented as a departure. The alternative is the definition given earlier: a system ends at its record, its key, its hardware and its channels.

Access is necessary and insufficient. A system with a dozen sources can still let a reinterpretation reach the commitment. Access removes only the case where the operator doesn't have to attack the boundary at all.
